% !TEX program = lualatex
% 自動生成: mar.report.latex — 手で編集した場合は再生成で上書きされる
\documentclass[paper=a4,fontsize=10.5pt,jafontsize=10.5pt]{jlreq}

\usepackage{luatexja}
\usepackage[haranoaji]{luatexja-preset}
\usepackage{amsmath,amssymb,amsthm,mathtools}
\usepackage{bm}
\usepackage{graphicx}
\usepackage{booktabs}
\usepackage{listings}
\usepackage{xcolor}
\usepackage[margin=25mm]{geometry}
\usepackage[hidelinks]{hyperref}
\usepackage{url}

\lstset{
  basicstyle=\ttfamily\footnotesize,
  breaklines=true,
  frame=single,
  showstringspaces=false,
  numbers=left,
  numberstyle=\tiny\color{gray},
  captionpos=b,
}

\theoremstyle{plain}
\newtheorem{theorem}{定理}[section]
\newtheorem{proposition}[theorem]{命題}
\newtheorem{lemma}[theorem]{補題}
\newtheorem{corollary}[theorem]{系}
\newtheorem{conjecture}[theorem]{予想}
\theoremstyle{definition}
\newtheorem{definition}[theorem]{定義}
\newtheorem{example}[theorem]{例}
\newtheorem{remark}[theorem]{注意}

\renewcommand{\proofname}{証明}

\title{最大誘導木の位数の下界: Written on the Wall II 予想 72・76・84・85 の部分的な証明と証人付き網羅検証}
\author{自動研究パイプライン \texttt{mar}（Claude Opus 5 + 検証器）}
\date{2026年07月28日}

\begin{document}
\maketitle

\begin{abstract}
DeLaViña の Graffiti.pc が生成した予想 Written on the Wall II Conjecture 72・76・84・85 を扱う (2026 年 7 月 27 日時点でいずれも未解決)。4 本とも最大誘導木の位数 $\mathrm{tree}(G)$ を距離量・次数量・局所独立数で下から抑える主張である。\textbf{まず手計算で 3 つの場合を落とす}。誘導測地路から出る $\mathrm{tree}(G) \ge \mathrm{diam}(G) + 1$ と、局所独立集合が張る誘導星から出る $\mathrm{tree}(G) \ge \ell_{\max} + 1$ を組み合わせると、予想 72 の弱い読みは\textbf{すべての連結グラフで成り立ち、しかも決して鋭くならない}こと、強い読みは $\ell_{\max} \le 3$ または $3\,\mathrm{diam}(G) \le 2\ell_{\max} + 3$ で成り立つこと、予想 84 は $\delta \ge 2$ で狭義に、$\delta = 1$ でも $\max(\mathrm{diam}, \ell_{\max}) \ge 2\,\mathrm{rad} - 1$ なら成り立つこと、予想 76 は木でつねに等号になることが従う。残るのは 72 が 1,487,419 個 (11.1\%)、84 が 15,080 個 (0.1\%) で、ほかに 76・85 である。これを、連結グラフの完全リスト ($n \le 10$)、木の完全リスト ($n \le 20$)、GENREG の連結正則グラフ、計 13,402,242 個に対して網羅検証した。反例は存在しない。$\mathrm{tree}(G)$ の計算は NP 困難だが、4 本の左辺はいずれも多項式時間で計算でき右辺は共通なので、\textbf{木を誘導する頂点集合 1 つ}を証人にすれば 4 本すべてが線形時間で確認できる。さらに下界が $\mathrm{tree}(G)$ にちょうど一致する場合を分類し、完全グラフ $K_n$ が予想 72 の強い読みと予想 85 の、偶数個の頂点をもつ道 $P_{2k}$ が予想 84 の、木が予想 76 の等号を与える無限族であることを示す。
\end{abstract}

\noindent\textbf{キーワード:} graph theory、induced tree、eccentricity、local independence、minimum degree、Graffiti.pc、open problem、certificate


\section{はじめに}

$G$ を位数 $n = |V(G)| \ge 2$ の連結な単純グラフとする。頂点部分集合
$T \subseteq V(G)$ が\textbf{木を誘導する}とは、誘導部分グラフ $G[T]$ が
連結かつ閉路をもたないことをいい、そのような $T$ の最大濃度を
$\mathrm{tree}(G)$ と書く。最大誘導木は古くから研究されている量であり
\cite{ess1986}、$n$ に比べていくらでも小さくなり得る
($\mathrm{tree}(K_n) = 2$)。その計算は NP 困難である。したがって
$\mathrm{tree}(G)$ の下界を多項式時間で計算できる量だけで与える主張は
自明ではない。

本稿が扱うのは、Graffiti.pc が生成した未解決予想集 Written on the Wall II
\cite{wowii} のうち、$\mathrm{tree}(G)$ の下界を与える 4 本である。
記号は同サイトの定義ページに従う。$\mathrm{ecc}(v)$ を離心数、
$\mathrm{diam}(G)$, $\mathrm{rad}(G)$ を直径・半径、$\delta(G)$ を最小次数、
$\deg_{\mathrm{avg}}(G) = 2|E(G)|/n$ を平均次数とし、さらに

\begin{itemize}
\item $\ell(v)$: $v$ の近傍が誘導する部分グラフ $G[N(v)]$ の独立数
      (\textbf{局所独立数})。$\ell_{\max} = \max_v \ell(v)$ と書く。
\item $T(v)$: $v$ を含む三角形の個数。$T_{\max} = \max_v T(v)$ とし、
      $\mathrm{freq}[T_{\max}]$ を $T(v) = T_{\max}$ なる頂点の個数とする。
\item $\mathrm{dist}_{\mathrm{even}}(v)$: $v$ からの距離が偶数の頂点の個数
      ($v$ 自身を含む)。
\end{itemize}

\begin{conjecture}[Graffiti.pc; WOWII Conjecture 76 \cite{wowii}]\label{conj:76}
$\dfrac{\mathrm{freq}[T_{\max}]}{\lfloor \deg_{\mathrm{avg}}(G) \rfloor}
 \ \le \ \mathrm{tree}(G)$.
\end{conjecture}

\begin{conjecture}[Graffiti.pc; WOWII Conjecture 84 \cite{wowii}]\label{conj:84}
$\dfrac{2\,\mathrm{rad}(G)}{\delta(G)} \ \le \ \mathrm{tree}(G)$.
\end{conjecture}

\begin{conjecture}[Graffiti.pc; WOWII Conjecture 85 \cite{wowii}]\label{conj:85}
$\left\lceil \sqrt{1 + 2 \min_{v} \mathrm{dist}_{\mathrm{even}}(v)}
 \right\rceil \ \le \ \mathrm{tree}(G)$.
\end{conjecture}

予想 72 の原文は \texttt{CEIL[average of ecc(v) + maximum of l (v) /3]} で
あり、\texttt{/3} が全体に係るのか $\ell_{\max}$ だけに係るのかが
\textbf{曖昧}である。本稿では両方を別の主張として扱う。

\begin{conjecture}[WOWII Conjecture 72、弱い読み]\label{conj:72a}
$\left\lceil \dfrac{1}{3}\left(
   \dfrac{1}{n}\sum_{v} \mathrm{ecc}(v) + \ell_{\max} \right)
 \right\rceil \ \le \ \mathrm{tree}(G)$.
\end{conjecture}

\begin{conjecture}[WOWII Conjecture 72、強い読み]\label{conj:72b}
$\left\lceil \dfrac{1}{n}\sum_{v} \mathrm{ecc}(v)
   + \dfrac{\ell_{\max}}{3} \right\rceil \ \le \ \mathrm{tree}(G)$.
\end{conjecture}

予想 \ref{conj:72b} は左辺が大きいので予想 \ref{conj:72a} を含意する。
2026 年 7 月 27 日に取得した出題者の公開ページ \cite{wowii} では、
4 本とも状態 \texttt{O} (未解決) である。

\begin{remark}[$n = 1$ の除外]\label{rem:k1}
$K_1$ では $\lfloor \deg_{\mathrm{avg}} \rfloor = \delta = 0$ となり、
予想 \ref{conj:76} と \ref{conj:84} の右辺が定義されない。以下つねに
$n \ge 2$ とする。
\end{remark}

計算上は分数と天井を次のように扱う。予想 \ref{conj:76}, \ref{conj:84} は
天井を含まないので分母を払って
\[ \mathrm{freq}[T_{\max}] \ \le \ \mathrm{tree}(G)
   \cdot \lfloor \deg_{\mathrm{avg}}(G) \rfloor, \qquad
   2\,\mathrm{rad}(G) \ \le \ \mathrm{tree}(G) \cdot \delta(G) \]
とし、予想 \ref{conj:72a}, \ref{conj:72b}, \ref{conj:85} は天井を
\textbf{先に取って整数に落としてから}比較する。本稿で\textbf{等号}
(tight) と呼ぶのは、この\emph{原文どおりの右辺}が $\mathrm{tree}(G)$ に
ちょうど一致することである。天井を取る前の実数が $\mathrm{tree}(G)$ より
真に小さくても等号は起こり得る。

\section{証人つき検証という設計}\label{sec:design}

\begin{lemma}\label{lem:witness}
$T \subseteq V(G)$ が木を誘導し、予想 \ref{conj:76}, \ref{conj:84},
\ref{conj:85}, \ref{conj:72a}, \ref{conj:72b} の $\mathrm{tree}(G)$ を
$|T|$ に置き換えた 5 本の不等式をすべて満たすならば、$G$ はいずれの予想の
反例でもない。
\end{lemma}

\begin{proof}
定義より $\mathrm{tree}(G) \ge |T|$。5 本はいずれも $\mathrm{tree}(G)$ に
ついて単調である ($\mathrm{tree}(G)$ は不等式の大きい側にしか現れず、
掛かる係数 $\lfloor \deg_{\mathrm{avg}} \rfloor$, $\delta$ は非負) から、
$|T|$ で成立すれば $\mathrm{tree}(G)$ でも成立する。
\end{proof}

仮定の確認に要するのは $G[T]$ が木かどうかの判定と、離心数・最小次数・
三角形数・局所独立数・偶数距離頂点数の計算である。局所独立数は誘導部分グラフ
$G[N(v)]$ の独立数なので一般には NP 困難な問題だが、走査した範囲では
$|N(v)| \le n-1$ の小さいグラフに対する計算であり支配的な費用にならない。
いずれにせよ、NP 困難な $\mathrm{tree}(G)$ そのものを検証器が解き直す必要は
ない。そこで本稿の証明書は、各グラフに対する $T$ を与える。証人は族ごとに
1 グラフ 4 バイト (little-endian の 32 ビット整数 1 個 = $T$ のビットマスク)
で列挙順に並べ、gzip 圧縮したバイナリに書き出す。そのファイルの SHA-256 を
証明書 JSON に記録し、検証器は同じ順序で元データを走査しながら証人を消費する。

\subsection{証人の作り方: 貪欲・木の近道・厳密解}

誘導木は「ちょうど 1 本の辺で現在の木につながる頂点」を足すことでしか育たない。
探索器は各頂点を根としてこの規則で貪欲に頂点を追加する。5 本すべてを
\emph{狭義}で満たすのに十分な大きさ $\mathrm{need}$ (各予想の右辺を
$\mathrm{tree}$ について解いた値の最大 $+\,1$) に届いた時点で打ち切ってよい。
そのグラフは反例でも等号でもないことが確定するからである。

貪欲が $\mathrm{need}$ に届かないグラフでは $\mathrm{tree}(G)$ を厳密に
求める必要があるが、\textbf{$G$ が木のときは全数探索が要らない}。
連結かつ $|E(G)| = n - 1$ なら $G$ 自身が木なので $V(G)$ が誘導木を与え、
$\mathrm{tree}(G) = n$ である。命題 \ref{prop:76tree} が示すように木は
必ず予想 \ref{conj:76} の等号になるので $\mathrm{need}$ には決して届かず、
この近道が無ければ木の族すべてで厳密解を呼ぶことになる。実際、走査した
13,402,242 個のうち近道が効いたのが 1,346,023 個、厳密解に進んだのは
16,254 個 (0.121\%) だけである。

\subsection{検証器}

検証器 (\texttt{mar.checkgraph}) は探索器のコードを一切参照せず、標準ライブラリ
のみで書かれている。アルゴリズムも意図的に変えてある。

\begin{enumerate}
\item 距離は Floyd--Warshall の全点対距離 (探索器はビットマスクの幅優先探索)。
      偶数距離の頂点数もこの距離行列から数える。
\item 三角形数は近傍の頂点対の総当たり (探索器はビット演算による共通近傍の
      ポップカウント)。
\item $\lceil a/b \rceil$ は商と剰余から組み立て、$\lceil \sqrt{s} \rceil$ は
      $r^2 \ge s$ となる最小の $r$ を単調に探す (探索器は整数平方根を使う)。
\item 木の判定は「森であり辺数 $= |T| - 1$」(探索器は「連結成分数 $= 1$ かつ森」)。
\item 最大誘導木は頂点ごとの採否による分枝 (探索器は葉を 1 枚ずつ生やす分枝)。
\end{enumerate}

元リストの個数を照合するための公表値 (OEIS A001349, A000055, A002851,
A006820--A006822) も、探索器の表ではなく検証器が独自にもつ定数と突き合わせる。
第 \ref{sec:hand} 節で証明する 4 つの主張についても、探索器と検証器の双方が
走査した全グラフで照合する (破れが 1 件でもあれば検証は失敗する)。

\subsection{等号の分類をどう閉じるか}\label{sec:tight}

補題 \ref{lem:witness} が閉じるのは不等式だけである。等号の主張は
$\mathrm{tree}(G)$ の上からの評価を含むので証人からは従わない。検証器は
\textbf{証人が 5 本すべてを狭義で満たさないグラフ}について
$\mathrm{tree}(G)$ を独立に計算し直し、予想ごとに狭義・等号・反例の 3 分類を
作って、その\textbf{件数が証明書の記録と一致すること}を確かめる。等号の
グラフを 1 個でも隠せば、そのグラフで狭義の不等式が破れるため件数が合わずに
検出される。加えて、族ごとの等号グラフの graph6 一覧が証明書にあるときは、
どの予想で等号になるかまで一致することを確かめる。tight なグラフがいずれかの予想で 2,000 個を超えた族 (8 族) では graph6 の一覧を省いた。そこでも分類は件数の照合で閉じる。

\section{手計算でどこまで詰められるか}\label{sec:hand}

網羅検証に入る前に、4 予想がどの範囲で\emph{証明できる}かを見る。
以下 $G$ は位数 $n \ge 2$ の連結グラフ、$t = \mathrm{tree}(G)$,
$D = \mathrm{diam}(G)$, $r = \mathrm{rad}(G)$, $\delta = \delta(G)$ とする。

\begin{lemma}\label{lem:diam}
$t \ge D + 1$.
\end{lemma}

\begin{proof}
最短路は誘導道である (途中の非隣接であるべき 2 頂点が隣接していれば、より短い
路が取れる)。距離 $D$ の 2 頂点を結ぶ最短路は $D+1$ 個の頂点からなる誘導道、
すなわち誘導木である。
\end{proof}

\begin{lemma}\label{lem:ell}
$t \ge \ell_{\max} + 1$.
\end{lemma}

\begin{proof}
$\ell(v) = \ell_{\max}$ となる $v$ を取り、$G[N(v)]$ の最大独立集合を $S$ と
する。$|S| = \ell_{\max}$ であり、$S \cup \{v\}$ が誘導する部分グラフは
$S$ の内部に辺が無く $v$ が $S$ の全点に隣接するから、星
$K_{1,\ell_{\max}}$ すなわち木である。
\end{proof}

\begin{theorem}\label{thm:72a}
予想 \ref{conj:72a} はすべての連結グラフ ($n \ge 2$) で成り立ち、しかも
\textbf{決して等号にならない}。
\end{theorem}

\begin{proof}
すべての $v$ で $\mathrm{ecc}(v) \le D$ だから平均離心数も $D$ 以下であり、
補題 \ref{lem:diam}, \ref{lem:ell} より $D \le t-1$ かつ
$\ell_{\max} \le t-1$。よって天井を取る前の値は
\[ \frac{1}{3}\left( \frac{1}{n}\sum_{v} \mathrm{ecc}(v)
   + \ell_{\max} \right) \ \le \ \frac{(t-1) + (t-1)}{3}
   \ = \ \frac{2t-2}{3} \ \le \ t - 1 \]
となる (最後の不等号は $t \ge 1$ と同値)。整数 $t-1$ 以下の実数の天井は
$t-1$ 以下だから、予想 \ref{conj:72a} の左辺は $t-1 < t$ であり、
狭義の不等式が成り立つ。
\end{proof}

\begin{theorem}\label{thm:72b}
$\ell_{\max} \le 3$ または $3D \le 2\ell_{\max} + 3$ ならば予想
\ref{conj:72b} が成り立つ。とくに反例 $G$ は $\ell_{\max} \ge 4$ かつ
$D \ge 4$、したがって $t \ge 5$ を満たす。
\end{theorem}

\begin{proof}
すべての $v$ で $\mathrm{ecc}(v) \le D$ だから平均離心数も $D$ 以下である。
$\ell_{\max} \le 3$ のときは補題 \ref{lem:diam} から
\[ \frac{1}{n}\sum_{v} \mathrm{ecc}(v) + \frac{\ell_{\max}}{3}
   \ \le \ (t-1) + \frac{3}{3} \ = \ t . \]
$3D \le 2\ell_{\max} + 3$ のときは $D \le \frac{2}{3}\ell_{\max} + 1$
なので $D + \ell_{\max}/3 \le \ell_{\max} + 1$ であり、補題
\ref{lem:ell} から
\[ \frac{1}{n}\sum_{v} \mathrm{ecc}(v) + \frac{\ell_{\max}}{3}
   \ \le \ D + \frac{\ell_{\max}}{3} \ \le \ \ell_{\max} + 1
   \ \le \ t . \]
どちらも右端が整数なので、天井を取っても $t$ を超えない。後半は対偶である。
$\ell_{\max} \ge 4$ と $3D > 2\ell_{\max} + 3 \ge 11$ から $D \ge 4$ が
出て、補題 \ref{lem:diam} より $t \ge 5$ となる。
\end{proof}

\begin{remark}[2 つの補題は合成できない]\label{rem:l3}
定理 \ref{thm:72b} が 2 つの枝に分かれるのは、補題 \ref{lem:diam} と
\ref{lem:ell} から得られる下界が $t \ge \max(D, \ell_{\max}) + 1$ 止まり
だからである。両者を足した $t \ge D + \ell_{\max} - 1$ が言えれば、
$\ell_{\max} \ge 2$ に対して
$\frac{1}{n}\sum_v \mathrm{ecc}(v) + \ell_{\max}/3
\le D + \ell_{\max}/3 \le t + 1 - 2\ell_{\max}/3 \le t$ となり、予想
\ref{conj:72b} は完全に片付く。しかしこの下界は\textbf{偽}である。
$4$ 閉路 $v_1 v_2 v_3 v_4$ の $v_4$ に葉 $v_5$ を 1 枚つけた $5$ 頂点グラフ
(graph6 表記 \texttt{DEw}) では、$\mathrm{dist}(v_2, v_5) = 3$ より $D = 3$、
$N(v_4) = \{v_1, v_3, v_5\}$ が独立集合だから $\ell_{\max} = 3$ である
一方、$G$ は閉路をもつので $t = 4 < 3 + 3 - 1$ となる。誘導測地路と誘導星は
一般には\textbf{重ならずに合成できない}。
\end{remark}

\begin{theorem}\label{thm:84}
$\delta \ge 2$ ならば予想 \ref{conj:84} は狭義に成り立つ。
\end{theorem}

\begin{proof}
$r \le D \le t - 1$ (補題 \ref{lem:diam}) と $\delta \ge 2$ より
\[ 2\,\mathrm{rad}(G) \ \le \ \delta \cdot r \ \le \ \delta (t-1)
   \ < \ \delta t . \]
\end{proof}

\begin{theorem}\label{thm:84d1}
$\delta = 1$ の場合も、$D \ge 2r - 1$ または $\ell_{\max} \ge 2r - 1$ なら
予想 \ref{conj:84} が成り立つ。とくに反例 $G$ は $\delta = 1$ かつ
$\max(D, \ell_{\max}) \le 2r - 2$ を満たす。
\end{theorem}

\begin{proof}
$\delta = 1$ なので右辺は $t$ である。$D \ge 2r - 1$ なら補題
\ref{lem:diam} より $t \ge D + 1 \ge 2r$、$\ell_{\max} \ge 2r - 1$ なら
補題 \ref{lem:ell} より $t \ge \ell_{\max} + 1 \ge 2r$ となる。
\end{proof}

\begin{remark}[残余は狭義の議論では閉じない]\label{rem:84resid}
定理 \ref{thm:84d1} の残余は空ではない。$6$ 閉路の 1 頂点に葉を 1 枚つけた
$7$ 頂点グラフ (graph6 表記 \texttt{F?qb\_}) は $\delta = 1$, $r = 3$, $D = 4$,
$\ell_{\max} = 3$ で、$D$ も $\ell_{\max}$ も $2r - 1 = 5$ に届かない。
しかもこのグラフは $t = 6 = 2r$ であり、予想 \ref{conj:84} の\textbf{等号}を
達成する。残余ケースが等号を達成するグラフを含む以上、狭義の不等式しか生まない
議論ではこの残余は閉じない。
\end{remark}

\begin{proposition}\label{prop:76tree}
$G$ が木 ($n \ge 2$) ならば予想 \ref{conj:76} は等号である。
\end{proposition}

\begin{proof}
木は三角形をもたないので全頂点で $T(v) = 0$、したがって $T_{\max} = 0$ かつ
$\mathrm{freq}[T_{\max}] = n$。また
$\deg_{\mathrm{avg}} = 2(n-1)/n = 2 - 2/n \in [1, 2)$ より
$\lfloor \deg_{\mathrm{avg}} \rfloor = 1$。$G$ 自身が誘導木だから $t = n$ で
あり、右辺は $n/1 = n = t$ となる。
\end{proof}

定理 \ref{thm:72a}, \ref{thm:72b}, \ref{thm:84}, \ref{thm:84d1} と命題
\ref{prop:76tree} は、走査した全グラフでも機械的に照合してある
(第 \ref{sec:design} 節)。未解決として残るのは
\textbf{$\ell_{\max} \ge 4$ かつ $3D > 2\ell_{\max} + 3$ の予想
\ref{conj:72b}}、\textbf{$\delta = 1$ かつ
$\max(D, \ell_{\max}) \le 2r - 2$ の予想 \ref{conj:84}}、および予想
\ref{conj:76}, \ref{conj:85} の一般の場合である。走査した 13,402,242 個の
うち、前者に該当するのが 1,487,419 個 (11.1\%)、後者が 15,080 個
(0.1\%) であった。

\section{網羅検証}

\begin{theorem}\label{thm:main}
位数 $n \le 10$ の連結グラフ全体 (11,989,763 個)、
位数 $11 \le n \le 20$ の木全体 (1,345,823 個。
位数 $10$ 以下の木は前者に含まれる)、
および $(n, r) \in \{(12, 3), (14, 3), (16, 3), (18, 3), (11, 4), (12, 4), (13, 4), (12, 5), (11, 6)\}$ の連結 $r$-正則グラフ全体
(66,656 個) のいずれにおいても、予想 \ref{conj:76}, \ref{conj:84},
\ref{conj:85} と予想 72 の 2 通りの読み (\ref{conj:72a},
\ref{conj:72b}) はすべて成立する。とくにこれらの反例 $G$ が存在するならば
$n \ge 11$ であり、$G$ が木ならば $n \ge 21$ である。
\end{theorem}

\begin{proof}
表 \ref{tab:fams} の各族について、木を誘導する頂点集合 $T$ を計算して記録した。
補題 \ref{lem:witness} より、各グラフで $G[T]$ が木であり $|T|$ が 5 本の
不等式を満たすことを確認すれば主張が従う。この確認は証明書
\texttt{p0011\_wowii\_tree\_lower\_bounds.json} と証人ファイルに対して検証器が実行し、
全 13,402,242 個で成立した。元リストの完全性は、走査行数が
OEIS A001349 (連結グラフ)、A000055 (木)、A002851 および
A006820--A006822 (連結正則グラフ) の公表値と一致することで担保される。
\end{proof}

\begin{table}[htbp]
\centering
\small
\caption{走査した族と、予想ごとに下界が $\mathrm{tree}(G)$ とちょうど一致した
グラフの個数。「厳密」は貪欲な証人が $\mathrm{need}$ に届かず最大誘導木を
厳密に解いた回数 (木の近道が効いた分は含まない)。}
\label{tab:fams}
\begin{tabular}{lrrrrrrrr}
\toprule
族 & $n$ & 個数 & 72a & 72b & 76 & 84 & 85 & 厳密 \\
\midrule
連結グラフ & 2 & 1 & 0 & 1 & 1 & 1 & 1 & 0 \\
連結グラフ & 3 & 2 & 0 & 2 & 1 & 0 & 1 & 1 \\
連結グラフ & 4 & 6 & 0 & 5 & 2 & 1 & 2 & 4 \\
連結グラフ & 5 & 21 & 0 & 11 & 3 & 4 & 2 & 11 \\
連結グラフ & 6 & 112 & 0 & 27 & 6 & 13 & 5 & 30 \\
連結グラフ & 7 & 853 & 0 & 86 & 11 & 40 & 6 & 102 \\
連結グラフ & 8 & 11,117 & 0 & 305 & 23 & 146 & 13 & 433 \\
連結グラフ & 9 & 261,080 & 0 & 970 & 47 & 561 & 17 & 2,022 \\
連結グラフ & 10 & 11,716,571 & 0 & 3,761 & 106 & 2,411 & 35 & 13,641 \\
木 & 11 & 235 & 0 & 0 & 235 & 0 & 0 & 0 \\
木 & 12 & 551 & 0 & 0 & 551 & 1 & 0 & 0 \\
木 & 13 & 1,301 & 0 & 0 & 1,301 & 0 & 0 & 0 \\
木 & 14 & 3,159 & 0 & 0 & 3,159 & 1 & 0 & 0 \\
木 & 15 & 7,741 & 0 & 0 & 7,741 & 0 & 0 & 0 \\
木 & 16 & 19,320 & 0 & 0 & 19,320 & 1 & 0 & 0 \\
木 & 17 & 48,629 & 0 & 0 & 48,629 & 0 & 0 & 0 \\
木 & 18 & 123,867 & 0 & 0 & 123,867 & 1 & 0 & 0 \\
木 & 19 & 317,955 & 0 & 0 & 317,955 & 0 & 0 & 0 \\
木 & 20 & 823,065 & 0 & 0 & 823,065 & 1 & 0 & 0 \\
3-正則 & 12 & 85 & 0 & 0 & 0 & 0 & 0 & 0 \\
3-正則 & 14 & 509 & 0 & 0 & 0 & 0 & 0 & 0 \\
3-正則 & 16 & 4,060 & 0 & 0 & 0 & 0 & 0 & 0 \\
3-正則 & 18 & 41,301 & 0 & 0 & 0 & 0 & 0 & 0 \\
4-正則 & 11 & 265 & 0 & 1 & 0 & 0 & 0 & 1 \\
4-正則 & 12 & 1,544 & 0 & 0 & 0 & 0 & 0 & 0 \\
4-正則 & 13 & 10,778 & 0 & 0 & 0 & 0 & 0 & 0 \\
5-正則 & 12 & 7,848 & 0 & 1 & 0 & 0 & 0 & 1 \\
6-正則 & 11 & 266 & 0 & 0 & 0 & 0 & 2 & 8 \\
\midrule
\multicolumn{2}{l}{合計} & 13,402,242 & 0
& 5,170 & 1,346,023 & 3,182 & 84
& 16,254 \\
\bottomrule
\end{tabular}
\end{table}

定理 \ref{thm:72a} が予言するとおり、走査した 13,402,242 個のいずれでも予想 \ref{conj:72a} は tight にならなかった (表 \ref{tab:fams} の 72a 列)。

命題 \ref{prop:76tree} が予言するとおり、木の族では 1,345,823 個すべてが予想 \ref{conj:76} の tight である (表 \ref{tab:fams})。

tight なグラフの例を graph6 表記で挙げる: 予想 \ref{conj:72b} では位数 $10$ の \texttt{I???F?]v?}、\texttt{I???F?]vO}、\texttt{I??CB@[f\_}、予想 \ref{conj:84} では位数 $10$ の \texttt{I???F?]vO}、\texttt{I??CB@[f?}、\texttt{I??CB@\textbackslash{}f\_}、予想 \ref{conj:85} では位数 $10$ の \texttt{I?AEDERaw}、\texttt{IQhTQj\textasciitilde{}\textasciitilde{}o}、\texttt{IQhTTV\textasciitilde{}\textasciitilde{}o} である。

\section{鋭さ: 等号を達成する無限族}\label{sec:sharp}

網羅検証は有限範囲の話だが、等号については無限族を手で与えられる。以下により、
予想 \ref{conj:72b}, \ref{conj:76}, \ref{conj:84}, \ref{conj:85} の
4 本すべてが\textbf{いくらでも大きい位数で改善不能}であることが分かる。

\begin{proposition}\label{prop:kn}
$n \ge 2$ のとき完全グラフ $K_n$ は予想 \ref{conj:72b} と \ref{conj:85} の
等号を達成する。
\end{proposition}

\begin{proof}
$K_n$ では全頂点の離心数が $1$ なので平均離心数は $1$、
$G[N(v)] = K_{n-1}$ より $\ell_{\max} = 1$、また $v$ 以外のすべての頂点が
$v$ から距離 $1$ にあるので $\mathrm{dist}_{\mathrm{even}}(v) = 1$ である。
誘導部分グラフが木になるのは高々 2 頂点なので $t = 2$。予想 \ref{conj:72b} の
左辺は $\lceil 1 + 1/3 \rceil = 2 = t$、予想 \ref{conj:85} の左辺は
$\lceil \sqrt{3} \rceil = 2 = t$ で、いずれも等号である。
\end{proof}

\begin{proposition}\label{prop:path}
$k \ge 1$ のとき $2k$ 個の頂点をもつ道 $P_{2k}$ は予想 \ref{conj:84} の
等号を達成する。奇数個の頂点をもつ道は達成しない。
\end{proposition}

\begin{proof}
$n$ 頂点の道の半径は $\lceil (n-1)/2 \rceil$、最小次数は $1$、そして道自身が
木なので $t = n$ である。$n = 2k$ なら
$2\,\mathrm{rad} = 2k = t \cdot \delta$ で等号、$n = 2k+1$ なら
$2\,\mathrm{rad} = 2k < 2k+1 = t \cdot \delta$ で狭義である。
\end{proof}

命題 \ref{prop:76tree} は、木という無限族が予想 \ref{conj:76} の等号を
達成することを既に与えている。一方、予想 \ref{conj:72a} については
定理 \ref{thm:72a} により等号を達成するグラフが\textbf{存在しない}。
すなわち予想 72 の 2 通りの読みのうち、\textbf{弱い読みは常に真だが決して
鋭くなく、強い読みは $K_n$ で鋭い}。出題者の意図が強い読みであることの
状況証拠と読める。

\section{限界}

本稿が示したのは第 \ref{sec:hand} 節の部分的な証明と、有限範囲での検証で
あって、4 予想の完全な証明ではない。連結グラフの完全リストは McKay
\cite{mckay} が $n \le 10$ まで、木は $n \le 22$ まで公開しており、
本稿は前者の全体と後者の $n \le 20$ を走査した。正則グラフは GENREG
\cite{genreg} が公開している範囲から、連結グラフの族に含まれない
$n \ge 11$ のものを選んだ。

一般の証明に必要な道具は予想ごとに異なる。予想 \ref{conj:72b} で残るのは
$\ell_{\max} \ge 4$ の場合であり、注意 \ref{rem:l3} が示すように
「誘導測地路と誘導星を合成する」という素朴な方針は使えない。必要なのは
$D$ と $\ell_{\max}$ が\emph{同時に}大きいときに働く下界である。
予想 \ref{conj:84} で残るのは $\delta = 1$、すなわち葉をもつグラフで
$t \ge 2\,\mathrm{rad}(G)$ を示すことであり、命題 \ref{prop:path} の道が
示すようにこの形は鋭い。予想 \ref{conj:76} と \ref{conj:85} については、
本稿は部分的な証明を与えられていない。

Graffiti.pc 自身がその開発時に小さいグラフのデータベース上で予想を試している
可能性が高い。したがって「位数の小さい連結グラフに反例がない」ことの新規性は
限定的であり、本稿の寄与は (i) 予想 72 の曖昧さを 2 通りの読みとして分離し、
弱い読みが\emph{常に真かつ決して鋭くない}ことを示したこと、(ii) 4 本を
1 つの証人で同時に検証できる形に落とし、第三者が独立に再実行できる証明書に
したこと、(iii) 部分的な証明により未解決部分を $\ell_{\max} \ge 4$ と
$\delta = 1$ に狭め、その残余が走査範囲でどれだけの割合を占めるかを測ったこと、
(iv) 等号を達成する無限族を予想 \ref{conj:72b}, \ref{conj:76},
\ref{conj:84}, \ref{conj:85} のすべてについて与えたこと、の 4 点にある。


\section*{再現性}
本稿の主張はすべて、探索器とは独立に実装された検証器によって再検査されている。次のコマンドで再現できる。

\begin{center}\texttt{python -m mar verify p0011\_wowii\_tree\_lower\_bounds}\end{center}

\begin{center}\begin{tabular}{ll}\toprule
証明書ダイジェスト & \texttt{3da23133ad0bba7a} \\
生成日時 (UTC) & 2026-07-27T16:35:00+00:00 \\
Python & 3.10.11 \\
実行環境 & Windows 11 (build 26200) \\
リポジトリ版 & \texttt{588a230} \\
乱数種 & 0 \\
探索時間 & 6290.0 秒 \\
\bottomrule\end{tabular}\end{center}


\begin{thebibliography}{99}
\bibitem{wowii} E. DeLaViña, Written on the Wall II: Conjectures of Graffiti.pc, University of Houston--Downtown (2026-07-27 取得). \url{http://cms.dt.uh.edu/faculty/delavinae/research/wowII/}
\bibitem{ess1986} P. Erdős, M. Saks, V. T. Sós, Maximum induced trees in graphs, J. Combin. Theory Ser. B 41 (1986) 61--79.
\bibitem{mckay} B. D. McKay, A. Piperno, Practical graph isomorphism II, J. Symbolic Comput. 60 (2014) 94--112. データ: Combinatorial Data. \url{https://users.cecs.anu.edu.au/~bdm/data/graphs.html}
\bibitem{genreg} M. Meringer, Fast generation of regular graphs and construction of cages, J. Graph Theory 30 (1999) 137--146. \url{https://www.mathe2.uni-bayreuth.de/markus/reggraphs.html}
\end{thebibliography}

\end{document}
